\cleardoubleoddpage
\begin{otherlanguage}{english}
\begin{abstract*}
Gyrokinetic plasma-turbulence simulations are computationally demanding. This thesis tests a
compact latent surrogate for eight-step forecasts of scalar flux and two one-dimensional spectra. A
JAX/Flax pipeline encodes preprocessed snapshots as 128-dimensional states; gated recurrent unit
(GRU) and causal Transformer models roll them forward through frozen diagnostic heads.

The study evaluates 51 trajectories by retrospective five-fold nested group cross-validation with
matched training seeds 52, 53, 54, 55, and 56; it is not a pristine locked-test result. Inner
validation selected Transformer in outer folds 0, 1, and 3 and GRU in folds 2 and 4. The selected
learned models achieved flux root mean squared error (RMSE) 14.6729, versus 10.0207 for copying the
last observed flux across the horizon. The difference was +4.6522 with paired hierarchical-bootstrap
interval [+2.5788, +6.6470]; learned error was lower in 31.6\% of hierarchically weighted paired
runs. The strictly positive interval supports a bounded negative result: the learned models do not
beat observed persistence.

Holding the last latent state fixed and decoding it gave flux RMSE 14.8049. Learned latent mean
squared error (MSE) was 0.4497 versus 0.4307 for latent persistence, a difference of +0.0189 with
interval [-0.0699, +0.1332]. Because this interval includes zero, no latent-dynamics advantage is
established. A diagnostic-head oracle given true future latents reached flux RMSE 12.3732, showing
usable diagnostic information but insufficient accuracy relative to observed persistence; it
isolates no single causal bottleneck.

The contribution is an auditable workflow, not a positive performance claim. Its scope is this
universe, preprocessing, and horizon; it establishes no full-field or long-horizon accuracy,
physical-unit or solver-level performance, broad generalization, a matched numerical comparison with
GyroSwin, or a universal architecture winner.
\end{abstract*}
\end{otherlanguage}

\cleardoubleoddpage
\begin{otherlanguage}{ngerman}
\begin{abstract*}
Gyrokinetische Simulationen von Plasmaturbulenz sind rechenintensiv. Diese Arbeit untersucht,
ob ein kompaktes latentes Ersatzmodell den skalaren Fluss und zwei eindimensionale Spektren über
acht Schritte vorhersagen kann. Eine JAX/Flax-Pipeline kodiert vorverarbeitete Momentaufnahmen als
128-dimensionale Zustände; GRU-Modelle (Gated Recurrent Units) und kausale Transformer-Modelle sagen
deren zeitliche Entwicklung voraus, und eingefrorene Diagnoseköpfe dekodieren die Zielgrößen.

Die Studie wertet 51 Trajektorien mittels retrospektiver fünffacher verschachtelter
Gruppen-Kreuzvalidierung mit denselben Trainings-Seeds 52, 53, 54, 55 und 56 aus; sie ist kein
Ergebnis einer bis zur Auswertung unangetasteten Testmenge. Die innere Validierung wählte den
Transformer in den äußeren Falten 0, 1 und 3 und das GRU-Modell in den Falten 2 und 4. Die
ausgewählten gelernten Modelle erreichten einen Fluss-RMSE (Wurzel des mittleren quadratischen
Fehlers) von 14.6729, während das Fortschreiben des letzten beobachteten Flusswerts über den
Horizont 10.0207 ergab. Die Differenz betrug +4.6522; das gepaarte hierarchische Bootstrap-Intervall
war [+2.5788, +6.6470]. In 31.6\% der hierarchisch gewichteten gepaarten Läufe war der Fehler des
gelernten Modells geringer. Das strikt positive Intervall stützt einen klar begrenzten negativen
Befund: Die gelernten Modelle übertreffen die beobachtete Persistenz nicht.

Das Festhalten am letzten latenten Zustand ergab nach der Dekodierung einen Fluss-RMSE von
14.8049. Der mittlere quadratische Fehler (MSE) im latenten Raum betrug 0.4497 für die gelernten
Modelle und 0.4307 für die latente Persistenz, entsprechend einer Differenz von +0.0189 mit dem
Intervall [-0.0699, +0.1332]. Da dieses Intervall Null einschließt, ist kein Vorteil der gelernten
latenten Dynamik belegt. Eine Oracle-Auswertung, bei der der Diagnosekopf die wahren zukünftigen
latenten Zustände erhielt, erreichte einen Fluss-RMSE von 12.3732. Dies zeigt, dass die Repräsentation
nutzbare diagnostische Informationen enthält, gegenüber der beobachteten Persistenz aber nicht genau
genug ist; einen einzelnen kausalen Engpass isoliert die Oracle-Auswertung nicht.

Der Beitrag liegt in einem überprüfbaren Arbeitsablauf, nicht in einem positiven Leistungsnachweis.
Seine Aussagekraft ist auf dieses Universum, diese Vorverarbeitung und diesen Horizont begrenzt; er
belegt weder vollständige Feld- oder Langzeitgenauigkeit noch Genauigkeit in physikalischen
Einheiten, Solver-Leistung, breite Verallgemeinerung, einen abgestimmten numerischen Vergleich mit
GyroSwin oder die allgemeingültige Überlegenheit einer Architektur.
\end{abstract*}
\end{otherlanguage}
